% --- Archon physics formalization source begin ---
% archon:physics
% archon:covers PhyXMiniProblems/problem_phyx_mini_0890.lean
% archon:source-report reports/phyx_mini/problem_phyx_mini_0890.source.json

\chapter{Physics problem phyx\_mini\_0890}
\label{ch:PhyXMiniProblems_problem_phyx_mini_0890}

\paragraph{Problem source.}
\#\# Physical scenario

Three charged particles with \textbackslash{}( q\_1 = -50\textbackslash{},\textbackslash{}mathrm\{nC\} \textbackslash{}), \textbackslash{}( q\_2 = +50\textbackslash{},\textbackslash{}mathrm\{nC\} \textbackslash{}), and \textbackslash{}( q\_3 = +30\textbackslash{},\textbackslash{}mathrm\{nC\} \textbackslash{}) are placed on the corners of the \textbackslash{}( 5.0\textbackslash{},\textbackslash{}mathrm\{cm\} \textbackslash{}times 10.0\textbackslash{},\textbackslash{}mathrm\{cm\} \textbackslash{}) rectangle shown in figure.

\#\# Question

What is the net force on charge \textbackslash{}( q\_3 \textbackslash{}) due to the other two charges?

\#\# Answer choices

- A: A: 1.25 \textbackslash{}times 10\textasciicircum{}\{-3\}N

- B: B: 1.35 \textbackslash{}times 10\textasciicircum{}\{-3\}N

- C: C: 6.2 \textbackslash{}times 10\textasciicircum{}\{-4\}N

- D: D: 1.33 \textbackslash{}times 10\textasciicircum{}\{-3\}N

\#\# Figure caption (auxiliary; use the image as primary evidence)

The image shows a diagram with three charges arranged in a right triangle. 

1. **Charges**:
   - Three point charges are labeled as \textbackslash{}( q\_1 \textbackslash{}), \textbackslash{}( q\_2 \textbackslash{}), and \textbackslash{}( q\_3 \textbackslash{}).
   - \textbackslash{}( q\_1 = -50 \textbackslash{}, \textbackslash{}text\{nC\} \textbackslash{}) (negative charge, represented in cyan).
   - \textbackslash{}( q\_2 = +50 \textbackslash{}, \textbackslash{}text\{nC\} \textbackslash{}) (positive charge, represented in red).
   - \textbackslash{}( q\_3 = +30 \textbackslash{}, \textbackslash{}text\{nC\} \textbackslash{}) (positive charge, represented in red).

2. **Distances**:
   - The horizontal distance between \textbackslash{}( q\_3 \textbackslash{}) and \textbackslash{}( q\_2 \textbackslash{}) is labeled as \textbackslash{}( 5.0 \textbackslash{}, \textbackslash{}text\{cm\} \textbackslash{}).
   - The vertical distance between \textbackslash{}( q\_1 \textbackslash{}) and \textbackslash{}( q\_3 \textbackslash{}) is labeled as \textbackslash{}( 10.0 \textbackslash{}, \textbackslash{}text\{cm\} \textbackslash{}).

3. **Diagram Layout**:
   - The charges form a right triangle with \textbackslash{}( q\_2 \textbackslash{}) and \textbackslash{}( q\_3 \textbackslash{}) aligned horizontally at the top, and \textbackslash{}( q\_1 \textbackslash{}) and \textbackslash{}( q\_3 \textbackslash{}) aligned vertically along the left side.
   - Dotted lines indicate the sides of the triangle, corresponding to the distances provided.

The image visually represents the arrangement of charges, useful for problems involving electrostatics or vector calculations.

\#\# Dataset metadata

Electrostatics; Spatial Relation Reasoning, Multi-Formula Reasoning

\paragraph{Recorded answer/context.}
C: C: 6.2 \textbackslash{}times 10\textasciicircum{}\{-4\}N

\paragraph{Figure/image path.}
/root/proposal\_for\_physic/science-mango/phyx\_mini\_run/phyx\_data/test\_image/890.png

\paragraph{Formalization target.}
create a compiling Lean file with sorry bodies at `PhyXMiniProblems/problem\_phyx\_mini\_0890.lean`. The Lean declarations must preserve the physical quantities, dimensions or dimensional roles, figure labels, governing-law hypotheses, and final relation expressed by this problem.
Use Mathlib/Physlib names found through LeanExplore where available. If a physics API is missing, introduce faithful local abstractions rather than scalar placeholder aliases.

\begin{theorem}[Physics formalization target]
\label{thm:physics:phyx_mini_0890:target}
\lean{PhyXMiniProblems.ProblemPhyXMini0890.problem_phyx_mini_0890}
\uses{def:physics:phyx-mini-0890:phyxminiproblems-problemphyxmini0890-forcevectorinnewtons, def:physics:phyx-mini-0890:phyxminiproblems-problemphyxmini0890-forcemagnitudeinnewtons, def:physics:phyx-mini-0890:phyxminiproblems-problemphyxmini0890-xaxis, def:physics:phyx-mini-0890:phyxminiproblems-problemphyxmini0890-yaxis, def:physics:phyx-mini-0890:phyxminiproblems-problemphyxmini0890-threechargerectanglesetup, def:physics:phyx-mini-0890:phyxminiproblems-problemphyxmini0890-hasthreechargerectangleproblemdata, def:physics:phyx-mini-0890:phyxminiproblems-problemphyxmini0890-matchesprimarythreechargefigure, def:physics:phyx-mini-0890:phyxminiproblems-problemphyxmini0890-satisfiesdisplayedrectanglegeometry, def:physics:phyx-mini-0890:phyxminiproblems-problemphyxmini0890-hasphysicalthreechargeparameters, def:physics:phyx-mini-0890:phyxminiproblems-problemphyxmini0890-usesroundedschoolcoulombconstant, def:physics:phyx-mini-0890:phyxminiproblems-problemphyxmini0890-satisfiespointchargecoulombandsuperpositionlaws, def:physics:phyx-mini-0890:phyxminiproblems-problemphyxmini0890-answerchoiceforceinnewtons, def:physics:phyx-mini-0890:phyxminiproblems-problemphyxmini0890-choicecdisplaytoleranceinnewtons, def:physics:phyx-mini-0890:phyxminiproblems-problemphyxmini0890-isclosestdisplayedforcemagnitude}
The assigned autoformalize agent should translate this physics problem into Lean declarations in the covered file, with theorem and lemma proof bodies written as `by sorry`.
\end{theorem}
\begin{proof}
This is an autoformalization task, not a proof task. Produce statements that can later be proved by the physics prover without weakening the physical model.
\end{proof}

% --- Archon declaration topology begin ---
\section{Declaration topology}
\begin{definition}[force Dimension]
  \label{def:physics:phyx-mini-0890:phyxminiproblems-problemphyxmini0890-forcedimension}
  \lean{PhyXMiniProblems.ProblemPhyXMini0890.forceDimension}
  The physical dimension M L T⁻² of force
\end{definition}
\begin{proof}
  This is the declared model definition of force Dimension.
\end{proof}
\begin{definition}[Signed Charge Quantity]
  \label{def:physics:phyx-mini-0890:phyxminiproblems-problemphyxmini0890-signedchargequantity}
  \lean{PhyXMiniProblems.ProblemPhyXMini0890.SignedChargeQuantity}
  A signed, unit-independent electric charge
\end{definition}
\begin{proof}
  This is the declared model definition of Signed Charge Quantity.
\end{proof}
\begin{definition}[Length Quantity]
  \label{def:physics:phyx-mini-0890:phyxminiproblems-problemphyxmini0890-lengthquantity}
  \lean{PhyXMiniProblems.ProblemPhyXMini0890.LengthQuantity}
  A nonnegative, unit-independent physical length
\end{definition}
\begin{proof}
  This is the declared model definition of Length Quantity.
\end{proof}
\begin{definition}[Planar Position Quantity]
  \label{def:physics:phyx-mini-0890:phyxminiproblems-problemphyxmini0890-planarpositionquantity}
  \lean{PhyXMiniProblems.ProblemPhyXMini0890.PlanarPositionQuantity}
  A unit-independent physical position in the plane of the figure
\end{definition}
\begin{proof}
  This is the declared model definition of Planar Position Quantity.
\end{proof}
\begin{definition}[Planar Force Quantity]
  \label{def:physics:phyx-mini-0890:phyxminiproblems-problemphyxmini0890-planarforcequantity}
  \lean{PhyXMiniProblems.ProblemPhyXMini0890.PlanarForceQuantity}
  \uses{def:physics:phyx-mini-0890:phyxminiproblems-problemphyxmini0890-forcedimension}
  A unit-independent physical force vector in the plane of the figure
\end{definition}
\begin{proof}
  This is the declared model definition of Planar Force Quantity.
\end{proof}
\begin{definition}[charge In Coulombs]
  \label{def:physics:phyx-mini-0890:phyxminiproblems-problemphyxmini0890-chargeincoulombs}
  \lean{PhyXMiniProblems.ProblemPhyXMini0890.chargeInCoulombs}
  \uses{def:physics:phyx-mini-0890:phyxminiproblems-problemphyxmini0890-signedchargequantity}
  Coherent-SI readout of a signed charge in coulombs
\end{definition}
\begin{proof}
  This is the declared model definition of charge In Coulombs.
\end{proof}
\begin{definition}[charge In Nanocoulombs]
  \label{def:physics:phyx-mini-0890:phyxminiproblems-problemphyxmini0890-chargeinnanocoulombs}
  \lean{PhyXMiniProblems.ProblemPhyXMini0890.chargeInNanocoulombs}
  \uses{def:physics:phyx-mini-0890:phyxminiproblems-problemphyxmini0890-signedchargequantity, def:physics:phyx-mini-0890:phyxminiproblems-problemphyxmini0890-chargeincoulombs}
  Nanocoulomb readout used by the three charge labels in the image
\end{definition}
\begin{proof}
  This is the declared model definition of charge In Nanocoulombs.
\end{proof}
\begin{definition}[length In Meters]
  \label{def:physics:phyx-mini-0890:phyxminiproblems-problemphyxmini0890-lengthinmeters}
  \lean{PhyXMiniProblems.ProblemPhyXMini0890.lengthInMeters}
  \uses{def:physics:phyx-mini-0890:phyxminiproblems-problemphyxmini0890-lengthquantity}
  Coherent-SI readout of a physical length in metres
\end{definition}
\begin{proof}
  This is the declared model definition of length In Meters.
\end{proof}
\begin{definition}[length In Centimeters]
  \label{def:physics:phyx-mini-0890:phyxminiproblems-problemphyxmini0890-lengthincentimeters}
  \lean{PhyXMiniProblems.ProblemPhyXMini0890.lengthInCentimeters}
  \uses{def:physics:phyx-mini-0890:phyxminiproblems-problemphyxmini0890-lengthquantity, def:physics:phyx-mini-0890:phyxminiproblems-problemphyxmini0890-lengthinmeters}
  Centimetre readout used by the two dimension labels in the image
\end{definition}
\begin{proof}
  This is the declared model definition of length In Centimeters.
\end{proof}
\begin{definition}[position Vector In Meters]
  \label{def:physics:phyx-mini-0890:phyxminiproblems-problemphyxmini0890-positionvectorinmeters}
  \lean{PhyXMiniProblems.ProblemPhyXMini0890.positionVectorInMeters}
  \uses{def:physics:phyx-mini-0890:phyxminiproblems-problemphyxmini0890-planarpositionquantity}
  Coherent-SI Cartesian position readout, whose coordinates are in metres
\end{definition}
\begin{proof}
  This is the declared model definition of position Vector In Meters.
\end{proof}
\begin{definition}[force Vector In Newtons]
  \label{def:physics:phyx-mini-0890:phyxminiproblems-problemphyxmini0890-forcevectorinnewtons}
  \lean{PhyXMiniProblems.ProblemPhyXMini0890.forceVectorInNewtons}
  \uses{def:physics:phyx-mini-0890:phyxminiproblems-problemphyxmini0890-planarforcequantity}
  Coherent-SI planar-force readout, whose coordinates are in newtons
\end{definition}
\begin{proof}
  This is the declared model definition of force Vector In Newtons.
\end{proof}
\begin{definition}[force Magnitude In Newtons]
  \label{def:physics:phyx-mini-0890:phyxminiproblems-problemphyxmini0890-forcemagnitudeinnewtons}
  \lean{PhyXMiniProblems.ProblemPhyXMini0890.forceMagnitudeInNewtons}
  \uses{def:physics:phyx-mini-0890:phyxminiproblems-problemphyxmini0890-planarforcequantity, def:physics:phyx-mini-0890:phyxminiproblems-problemphyxmini0890-forcevectorinnewtons}
  Euclidean magnitude in newtons of a dimensionful planar force
\end{definition}
\begin{proof}
  This is the declared model definition of force Magnitude In Newtons.
\end{proof}
\begin{definition}[x Axis]
  \label{def:physics:phyx-mini-0890:phyxminiproblems-problemphyxmini0890-xaxis}
  \lean{PhyXMiniProblems.ProblemPhyXMini0890.xAxis}
  Coordinate 0, horizontal and positive to the right
\end{definition}
\begin{proof}
  This is the declared model definition of x Axis.
\end{proof}
\begin{definition}[y Axis]
  \label{def:physics:phyx-mini-0890:phyxminiproblems-problemphyxmini0890-yaxis}
  \lean{PhyXMiniProblems.ProblemPhyXMini0890.yAxis}
  Coordinate 1, vertical and positive upward
\end{definition}
\begin{proof}
  This is the declared model definition of y Axis.
\end{proof}
\begin{definition}[Particle]
  \label{def:physics:phyx-mini-0890:phyxminiproblems-problemphyxmini0890-particle}
  \lean{PhyXMiniProblems.ProblemPhyXMini0890.Particle}
  The three charged particles named in the source and the image
\end{definition}
\begin{proof}
  This is the declared model definition of Particle.
\end{proof}
\begin{definition}[Rectangle Corner]
  \label{def:physics:phyx-mini-0890:phyxminiproblems-problemphyxmini0890-rectanglecorner}
  \lean{PhyXMiniProblems.ProblemPhyXMini0890.RectangleCorner}
  The four corners of the dashed rectangle
\end{definition}
\begin{proof}
  This is the declared model definition of Rectangle Corner.
\end{proof}
\begin{definition}[Rectangle Side]
  \label{def:physics:phyx-mini-0890:phyxminiproblems-problemphyxmini0890-rectangleside}
  \lean{PhyXMiniProblems.ProblemPhyXMini0890.RectangleSide}
  The four dashed sides visible in the primary image
\end{definition}
\begin{proof}
  This is the declared model definition of Rectangle Side.
\end{proof}
\begin{definition}[Figure Charge Glyph]
  \label{def:physics:phyx-mini-0890:phyxminiproblems-problemphyxmini0890-figurechargeglyph}
  \lean{PhyXMiniProblems.ProblemPhyXMini0890.FigureChargeGlyph}
  The plus or minus glyph drawn inside each charge circle
\end{definition}
\begin{proof}
  This is the declared model definition of Figure Charge Glyph.
\end{proof}
\begin{definition}[Charged Particle Model]
  \label{def:physics:phyx-mini-0890:phyxminiproblems-problemphyxmini0890-chargedparticlemodel}
  \lean{PhyXMiniProblems.ProblemPhyXMini0890.ChargedParticleModel}
  The physical idealization used when applying the point-charge law
\end{definition}
\begin{proof}
  This is the declared model definition of Charged Particle Model.
\end{proof}
\begin{definition}[expected Corner]
  \label{def:physics:phyx-mini-0890:phyxminiproblems-problemphyxmini0890-expectedcorner}
  \lean{PhyXMiniProblems.ProblemPhyXMini0890.expectedCorner}
  \uses{def:physics:phyx-mini-0890:phyxminiproblems-problemphyxmini0890-particle, def:physics:phyx-mini-0890:phyxminiproblems-problemphyxmini0890-rectanglecorner}
  Corner occupied by each named particle in image 890.png
\end{definition}
\begin{proof}
  This is the declared model definition of expected Corner.
\end{proof}
\begin{definition}[expected Charge Glyph]
  \label{def:physics:phyx-mini-0890:phyxminiproblems-problemphyxmini0890-expectedchargeglyph}
  \lean{PhyXMiniProblems.ProblemPhyXMini0890.expectedChargeGlyph}
  \uses{def:physics:phyx-mini-0890:phyxminiproblems-problemphyxmini0890-particle, def:physics:phyx-mini-0890:phyxminiproblems-problemphyxmini0890-figurechargeglyph}
  Charge-sign glyph visibly attached to each particle
\end{definition}
\begin{proof}
  This is the declared model definition of expected Charge Glyph.
\end{proof}
\begin{definition}[Three Charge Rectangle Figure]
  \label{def:physics:phyx-mini-0890:phyxminiproblems-problemphyxmini0890-threechargerectanglefigure}
  \lean{PhyXMiniProblems.ProblemPhyXMini0890.ThreeChargeRectangleFigure}
  \uses{def:physics:phyx-mini-0890:phyxminiproblems-problemphyxmini0890-particle, def:physics:phyx-mini-0890:phyxminiproblems-problemphyxmini0890-rectanglecorner, def:physics:phyx-mini-0890:phyxminiproblems-problemphyxmini0890-rectangleside, def:physics:phyx-mini-0890:phyxminiproblems-problemphyxmini0890-figurechargeglyph}
  Literal presentation content transcribed from the primary raster. The image does not display a net-force vector, magnitude, or answer choice
\end{definition}
\begin{proof}
  This is the declared model definition of Three Charge Rectangle Figure.
\end{proof}
\begin{definition}[Three Charge Rectangle Setup]
  \label{def:physics:phyx-mini-0890:phyxminiproblems-problemphyxmini0890-threechargerectanglesetup}
  \lean{PhyXMiniProblems.ProblemPhyXMini0890.ThreeChargeRectangleSetup}
  \uses{def:physics:phyx-mini-0890:phyxminiproblems-problemphyxmini0890-signedchargequantity, def:physics:phyx-mini-0890:phyxminiproblems-problemphyxmini0890-lengthquantity, def:physics:phyx-mini-0890:phyxminiproblems-problemphyxmini0890-planarpositionquantity, def:physics:phyx-mini-0890:phyxminiproblems-problemphyxmini0890-planarforcequantity, def:physics:phyx-mini-0890:phyxminiproblems-problemphyxmini0890-particle, def:physics:phyx-mini-0890:phyxminiproblems-problemphyxmini0890-chargedparticlemodel, def:physics:phyx-mini-0890:phyxminiproblems-problemphyxmini0890-threechargerectanglefigure}
  Independent physical objects in the electrostatic setup. In particular, netForceOn and pairForceOnDueTo are observables constrained by the laws below; neither is defined from a displayed answer
\end{definition}
\begin{proof}
  This is the declared model definition of Three Charge Rectangle Setup.
\end{proof}
\begin{definition}[displacement Vector In Meters]
  \label{def:physics:phyx-mini-0890:phyxminiproblems-problemphyxmini0890-displacementvectorinmeters}
  \lean{PhyXMiniProblems.ProblemPhyXMini0890.displacementVectorInMeters}
  \uses{def:physics:phyx-mini-0890:phyxminiproblems-problemphyxmini0890-positionvectorinmeters, def:physics:phyx-mini-0890:phyxminiproblems-problemphyxmini0890-particle, def:physics:phyx-mini-0890:phyxminiproblems-problemphyxmini0890-threechargerectanglesetup}
  Displacement from a source particle to a target particle, in metres
\end{definition}
\begin{proof}
  This is the declared model definition of displacement Vector In Meters.
\end{proof}
\begin{definition}[separation In Meters]
  \label{def:physics:phyx-mini-0890:phyxminiproblems-problemphyxmini0890-separationinmeters}
  \lean{PhyXMiniProblems.ProblemPhyXMini0890.separationInMeters}
  \uses{def:physics:phyx-mini-0890:phyxminiproblems-problemphyxmini0890-particle, def:physics:phyx-mini-0890:phyxminiproblems-problemphyxmini0890-threechargerectanglesetup, def:physics:phyx-mini-0890:phyxminiproblems-problemphyxmini0890-displacementvectorinmeters}
  Euclidean separation of two particles, in metres
\end{definition}
\begin{proof}
  This is the declared model definition of separation In Meters.
\end{proof}
\begin{definition}[Has Three Charge Rectangle Problem Data]
  \label{def:physics:phyx-mini-0890:phyxminiproblems-problemphyxmini0890-hasthreechargerectangleproblemdata}
  \lean{PhyXMiniProblems.ProblemPhyXMini0890.HasThreeChargeRectangleProblemData}
  \uses{def:physics:phyx-mini-0890:phyxminiproblems-problemphyxmini0890-chargeinnanocoulombs, def:physics:phyx-mini-0890:phyxminiproblems-problemphyxmini0890-lengthincentimeters, def:physics:phyx-mini-0890:phyxminiproblems-problemphyxmini0890-threechargerectanglesetup}
  The numerical physical data stated in the problem. These are charge and length readouts only and contain no force value
\end{definition}
\begin{proof}
  This is the declared model definition of Has Three Charge Rectangle Problem Data.
\end{proof}
\begin{definition}[Matches Primary Three Charge Figure]
  \label{def:physics:phyx-mini-0890:phyxminiproblems-problemphyxmini0890-matchesprimarythreechargefigure}
  \lean{PhyXMiniProblems.ProblemPhyXMini0890.MatchesPrimaryThreeChargeFigure}
  \uses{def:physics:phyx-mini-0890:phyxminiproblems-problemphyxmini0890-chargeinnanocoulombs, def:physics:phyx-mini-0890:phyxminiproblems-problemphyxmini0890-lengthincentimeters, def:physics:phyx-mini-0890:phyxminiproblems-problemphyxmini0890-expectedcorner, def:physics:phyx-mini-0890:phyxminiproblems-problemphyxmini0890-expectedchargeglyph, def:physics:phyx-mini-0890:phyxminiproblems-problemphyxmini0890-threechargerectanglesetup}
  Unambiguous evidence read from 890.png. In contrast with the auxiliary caption, the primary raster places q₂ at the bottom-right corner, so the line from q₂ to q₃ is the rectangle diagonal
\end{definition}
\begin{proof}
  This is the declared model definition of Matches Primary Three Charge Figure.
\end{proof}
\begin{definition}[Satisfies Displayed Rectangle Geometry]
  \label{def:physics:phyx-mini-0890:phyxminiproblems-problemphyxmini0890-satisfiesdisplayedrectanglegeometry}
  \lean{PhyXMiniProblems.ProblemPhyXMini0890.SatisfiesDisplayedRectangleGeometry}
  \uses{def:physics:phyx-mini-0890:phyxminiproblems-problemphyxmini0890-lengthinmeters, def:physics:phyx-mini-0890:phyxminiproblems-problemphyxmini0890-positionvectorinmeters, def:physics:phyx-mini-0890:phyxminiproblems-problemphyxmini0890-xaxis, def:physics:phyx-mini-0890:phyxminiproblems-problemphyxmini0890-yaxis, def:physics:phyx-mini-0890:phyxminiproblems-problemphyxmini0890-threechargerectanglesetup}
  The coordinate consequences of the occupied rectangle corners. Translation of the whole diagram remains free: only the relative horizontal and vertical coordinates are fixed
\end{definition}
\begin{proof}
  This is the declared model definition of Satisfies Displayed Rectangle Geometry.
\end{proof}
\begin{definition}[Has Physical Three Charge Parameters]
  \label{def:physics:phyx-mini-0890:phyxminiproblems-problemphyxmini0890-hasphysicalthreechargeparameters}
  \lean{PhyXMiniProblems.ProblemPhyXMini0890.HasPhysicalThreeChargeParameters}
  \uses{def:physics:phyx-mini-0890:phyxminiproblems-problemphyxmini0890-chargeincoulombs, def:physics:phyx-mini-0890:phyxminiproblems-problemphyxmini0890-lengthinmeters, def:physics:phyx-mini-0890:phyxminiproblems-problemphyxmini0890-threechargerectanglesetup, def:physics:phyx-mini-0890:phyxminiproblems-problemphyxmini0890-separationinmeters}
  Point-particle and nondegeneracy conditions selecting the physical branch
\end{definition}
\begin{proof}
  This is the declared model definition of Has Physical Three Charge Parameters.
\end{proof}
\begin{definition}[Uses Rounded School Coulomb Constant]
  \label{def:physics:phyx-mini-0890:phyxminiproblems-problemphyxmini0890-usesroundedschoolcoulombconstant}
  \lean{PhyXMiniProblems.ProblemPhyXMini0890.UsesRoundedSchoolCoulombConstant}
  \uses{def:physics:phyx-mini-0890:phyxminiproblems-problemphyxmini0890-threechargerectanglesetup}
  The rounded Coulomb constant customarily used for the multiple-choice calculation, expressed in coherent SI units N m² / C²
\end{definition}
\begin{proof}
  This is the declared model definition of Uses Rounded School Coulomb Constant.
\end{proof}
\begin{definition}[Satisfies Point Charge Coulomb And Superposition Laws]
  \label{def:physics:phyx-mini-0890:phyxminiproblems-problemphyxmini0890-satisfiespointchargecoulombandsuperpositionlaws}
  \lean{PhyXMiniProblems.ProblemPhyXMini0890.SatisfiesPointChargeCoulombAndSuperpositionLaws}
  \uses{def:physics:phyx-mini-0890:phyxminiproblems-problemphyxmini0890-chargeincoulombs, def:physics:phyx-mini-0890:phyxminiproblems-problemphyxmini0890-forcevectorinnewtons, def:physics:phyx-mini-0890:phyxminiproblems-problemphyxmini0890-particle, def:physics:phyx-mini-0890:phyxminiproblems-problemphyxmini0890-threechargerectanglesetup, def:physics:phyx-mini-0890:phyxminiproblems-problemphyxmini0890-displacementvectorinmeters, def:physics:phyx-mini-0890:phyxminiproblems-problemphyxmini0890-separationinmeters}
  For distinct point charges, the signed vector Coulomb law is F(target <- source) = k q\_target q\_source (r\_target - r\_source) / |r|\textasciicircum{}3. The second field states vector superposition over every other named particle. These are general governing laws and contain no requested numerical force or answer choice
\end{definition}
\begin{proof}
  This is the declared model definition of Satisfies Point Charge Coulomb And Superposition Laws.
\end{proof}
\begin{lemma}[q3 pair separations]
  \label{lem:physics:phyx-mini-0890:phyxminiproblems-problemphyxmini0890-q3-pair-separations}
  \lean{PhyXMiniProblems.ProblemPhyXMini0890.q3_pair_separations}
  \uses{def:physics:phyx-mini-0890:phyxminiproblems-problemphyxmini0890-lengthinmeters, def:physics:phyx-mini-0890:phyxminiproblems-problemphyxmini0890-threechargerectanglesetup, def:physics:phyx-mini-0890:phyxminiproblems-problemphyxmini0890-separationinmeters, def:physics:phyx-mini-0890:phyxminiproblems-problemphyxmini0890-satisfiesdisplayedrectanglegeometry, def:physics:phyx-mini-0890:phyxminiproblems-problemphyxmini0890-hasphysicalthreechargeparameters}
  The two source-to-target separations are respectively the rectangle height and its diagonal. This follows from the primary-image corner layout and is not an independent figure readout
\end{lemma}
\begin{proof}
  The conclusion follows from the governing relations and figure readouts named in the statement of q3 pair separations.
\end{proof}
\begin{lemma}[net Force On Q3 component formulas]
  \label{lem:physics:phyx-mini-0890:phyxminiproblems-problemphyxmini0890-netforceonq3-component-formulas}
  \lean{PhyXMiniProblems.ProblemPhyXMini0890.netForceOnQ3_component_formulas}
  \uses{def:physics:phyx-mini-0890:phyxminiproblems-problemphyxmini0890-chargeincoulombs, def:physics:phyx-mini-0890:phyxminiproblems-problemphyxmini0890-lengthinmeters, def:physics:phyx-mini-0890:phyxminiproblems-problemphyxmini0890-forcevectorinnewtons, def:physics:phyx-mini-0890:phyxminiproblems-problemphyxmini0890-xaxis, def:physics:phyx-mini-0890:phyxminiproblems-problemphyxmini0890-yaxis, def:physics:phyx-mini-0890:phyxminiproblems-problemphyxmini0890-threechargerectanglesetup, def:physics:phyx-mini-0890:phyxminiproblems-problemphyxmini0890-satisfiesdisplayedrectanglegeometry, def:physics:phyx-mini-0890:phyxminiproblems-problemphyxmini0890-hasphysicalthreechargeparameters, def:physics:phyx-mini-0890:phyxminiproblems-problemphyxmini0890-satisfiespointchargecoulombandsuperpositionlaws}
  Resolving the two Coulomb forces along the displayed axes gives these exact component formulas. The negative q₁ term is automatically attractive because its signed charge readout occurs in the formula
\end{lemma}
\begin{proof}
  The conclusion follows from the governing relations and figure readouts named in the statement of net Force On Q3 component formulas.
\end{proof}
\begin{definition}[Answer Choice]
  \label{def:physics:phyx-mini-0890:phyxminiproblems-problemphyxmini0890-answerchoice}
  \lean{PhyXMiniProblems.ProblemPhyXMini0890.AnswerChoice}
  Labels attached to the four displayed force-magnitude choices
\end{definition}
\begin{proof}
  This is the declared model definition of Answer Choice.
\end{proof}
\begin{definition}[answer Choice Force In Newtons]
  \label{def:physics:phyx-mini-0890:phyxminiproblems-problemphyxmini0890-answerchoiceforceinnewtons}
  \lean{PhyXMiniProblems.ProblemPhyXMini0890.answerChoiceForceInNewtons}
  \uses{def:physics:phyx-mini-0890:phyxminiproblems-problemphyxmini0890-answerchoice}
  Literal force magnitude in newtons printed beside each answer label
\end{definition}
\begin{proof}
  This is the declared model definition of answer Choice Force In Newtons.
\end{proof}
\begin{definition}[choice CDisplay Tolerance In Newtons]
  \label{def:physics:phyx-mini-0890:phyxminiproblems-problemphyxmini0890-choicecdisplaytoleranceinnewtons}
  \lean{PhyXMiniProblems.ProblemPhyXMini0890.choiceCDisplayToleranceInNewtons}
  Half of the last decimal place displayed in choice C
\end{definition}
\begin{proof}
  This is the declared model definition of choice CDisplay Tolerance In Newtons.
\end{proof}
\begin{definition}[Is Closest Displayed Force Magnitude]
  \label{def:physics:phyx-mini-0890:phyxminiproblems-problemphyxmini0890-isclosestdisplayedforcemagnitude}
  \lean{PhyXMiniProblems.ProblemPhyXMini0890.IsClosestDisplayedForceMagnitude}
  \uses{def:physics:phyx-mini-0890:phyxminiproblems-problemphyxmini0890-planarforcequantity, def:physics:phyx-mini-0890:phyxminiproblems-problemphyxmini0890-forcemagnitudeinnewtons, def:physics:phyx-mini-0890:phyxminiproblems-problemphyxmini0890-answerchoice, def:physics:phyx-mini-0890:phyxminiproblems-problemphyxmini0890-answerchoiceforceinnewtons}
  A choice is closest when its printed magnitude is no farther from the independently modeled net-force magnitude than every alternative
\end{definition}
\begin{proof}
  This is the declared model definition of Is Closest Displayed Force Magnitude.
\end{proof}
% --- Archon declaration topology end ---
% --- Archon physics formalization source end ---
